% !TEX program = xelatex % Указываем, что компилятор - XeLaTeX
\documentclass[aspectratio=169]{beamer}
\usetheme{Madrid}
\usecolortheme{beaver}

% ---------- Encoding & language (xelatex-friendly) ----------
\usepackage{polyglossia} % Вместо babel для XeLaTeX
\setdefaultlanguage{russian}
\setotherlanguages{english}

\usepackage{fontspec} % Для работы со шрифтами
\setmainfont{DejaVu Serif}
\setsansfont{DejaVu Sans}
\setmonofont{DejaVu Sans Mono}

% ---------- Common packages ----------
\usepackage{graphicx}
\usepackage{booktabs}
\usepackage{tabularx}
\usepackage{hyperref}
\usepackage{multicol}
\usepackage{ulem}
\usepackage{csquotes}
\usepackage{tikz}
\usetikzlibrary{arrows.meta,positioning,fit,shapes.misc}

% ---------- Listings (Kotlin/Java) ----------
\usepackage{listings} % Используем стандартный listings (inputenc не нужен)
\usepackage[dvipsnames]{xcolor}

\lstdefinelanguage{Kotlin}{
  comment=[l]{//},
  commentstyle={\color{gray}\ttfamily},
  emph={filter, first, firstOrNull, forEach, lazy, map, mapNotNull, println, launch, async, await, withContext},
  emphstyle={\color{OrangeRed}},
  identifierstyle=\color{black},
  keywords={!in, !is, abstract, actual, annotation, as, as?, break, by, catch, class, companion, const, constructor, continue, crossinline, data, delegate, do, dynamic, else, enum, expect, external, false, field, file, final, finally, for, fun, get, if, import, in, infix, init, inline, inner, interface, internal, is, lateinit, noinline, null, object, open, operator, out, override, package, param, private, property, protected, public, receiver, reified, return, return@, sealed, set, setparam, super, suspend, tailrec, this, throw, true, try, typealias, typeof, val, var, vararg, when, where, while},
  keywordstyle={\color{NavyBlue}\bfseries},
  morecomment=[s]{/*}{*/},
  morestring=[b]",
  morestring=[s]{"""*}{*"""},
  ndkeywords={@Composable, @Preview, @Deprecated, @JvmField, @JvmName, @JvmOverloads, @JvmStatic, @JvmSynthetic, Array, Byte, Double, Float, Int, Integer, Iterable, Long, Runnable, Short, String, Any, Unit, Nothing},
  ndkeywordstyle={\color{BurntOrange}\bfseries},
  sensitive=true,
  stringstyle={\color{ForestGreen}\ttfamily}
}

\lstdefinelanguage{JavaLite}{
  language=Java,
  basicstyle=\footnotesize\ttfamily,
  keywordstyle=\color{NavyBlue}\bfseries,
  commentstyle=\color{gray}\ttfamily,
  stringstyle=\color{ForestGreen}\ttfamily
}

\lstset{
  % inputencoding=utf8, % Не нужно для XeLaTeX
  extendedchars=true,
  basicstyle=\footnotesize\ttfamily,
  breaklines=true,
  showstringspaces=false,
  frame=single,
  framerule=0.2pt,
  tabsize=2,
  numbers=left,
  numberstyle=\tiny,
  numbersep=6pt
}

% ---------- Title ----------
\title{Введение в мобильную разработку (Android/Kotlin)}
\subtitle{Лекция 1: Подходы к созданию мобильных приложений, структура курса и обзор Kotlin}
\author{Александр Глускер}
\institute{Курс «Мобильная разработка»}
\date{\today}

% ---------- TOC at section start ----------
\AtBeginSection[]{
  \begin{frame}{Навигация по лекции}
    \tableofcontents[currentsection, hideallsubsections]
  \end{frame}
}

\begin{document}

% ================== Title ==================
\begin{frame}
  \titlepage
\end{frame}

% ================== Agenda ==================
\begin{frame}{Содержание}
  \begin{columns}[T,onlytextwidth]
    \begin{column}{0.55\textwidth}
      \begin{enumerate}
        \item Формальности и знакомство 
        \item Введение 
        \item Подходы к мобильной разработке
        \item Kotlin: обзор языка и примеры
        \item Итоги, Q\&A
      \end{enumerate}
    \end{column}
    \begin{column}{0.45\textwidth}
      \begin{block}{Цели лекции}
        \begin{itemize}
          \item Разобраться в требованиях курса и оценивании
          \item Понять ландшафт мобильной разработки
          \item Увидеть особенности Kotlin
        \end{itemize}
      \end{block}
    \end{column}
  \end{columns}
\end{frame}

\section{Формальности и знакомство}

\begin{frame}{Как мы учимся на курсе}
  \begin{itemize}
    \item \textbf{Лекции}: теория (только главное), обзор, мастер-классы (но на них будет не всё)
    \item \textbf{Семинары}: разбор заданий (если все будет хорошо), сдача
    \item \textbf{Дома}: выполнение задания
  \end{itemize}
  \begin{block}{Инструменты}
    IntelliJ Idea, Android Studio
  \end{block}
\end{frame}

\begin{frame}{Оценивание и контроль}
  \begin{itemize}
    \item \textbf{Практические задания:} на каждом семинаре
    \item \textbf{Полусеместр} письменный тест (от руки) + накопленная оценка за практику
    \item \textbf{Финальный экзамен:} или с нуля (теория + практика), или накопленная оценка
    \item \textbf{Проекты etc:} возможны к учету
  \end{itemize}
\end{frame}

\begin{frame}{Политика курса}
  \begin{itemize}
    \item Дедлайны по практикам, их можно сдавать позже, но оценка ниже и приоритет тем, кто сдает вовремя
    \item Код должен соответствовать стандартам (обсудим)
    \item Как можно использовать AI (методика проверки)
    \item Вопросы в начале пар
  \end{itemize}
\end{frame}


% =====================================================
\section{Подходы к мобильной разработке}
% =====================================================

\begin{frame}{Ландшафт: основные подходы}
  \begin{itemize}
    \item Нативная разработка (Android: Kotlin/Java; iOS: Swift)
    \item Кросс-платформенная (Flutter, React Native, Kotlin Multiplatform, .NET MAUI)
    \item Гибридная (Ionic/Cordova)
    \item PWA (продвинутые веб-приложения)
    \item Low-code/No-code
  \end{itemize}
\end{frame}

\begin{frame}{Нативная разработка: архитектура}
  \centering
  \begin{tikzpicture}[node distance=5mm, every node/.style={align=center}]
    \node[draw, rounded corners, fill=gray!10, minimum width=0.9\linewidth, inner sep=3mm] (app) {Приложение (Kotlin/Swift) \\ + UI (Jetpack Compose/SwiftUI) \\ + Домены и Use Cases};
    \node[draw, rounded corners, fill=gray!10, below=of app, minimum width=0.9\linewidth, inner sep=3mm] (sdk) {SDK/Фреймворки платформы (Android Jetpack / iOS Frameworks)};
    \node[draw, rounded corners, fill=gray!10, below=of sdk, minimum width=0.9\linewidth, inner sep=3mm] (os) {ОС и драйверы (Доступ к датчикам, сети, файловой системе)};
    \draw[-{Latex}] (app) -- (sdk);
    \draw[-{Latex}] (sdk) -- (os);
  \end{tikzpicture}
  \vspace{3mm}

  \begin{itemize}
    \item \textbf{Плюсы:} производительность, полный доступ к API, лучший UX
    \item \textbf{Минусы:} отдельные коды для iOS/Android, выше стоимость
  \end{itemize}
\end{frame}

\begin{frame}{Кросс-платформенные решения}
  \begin{columns}[T,onlytextwidth]
    \begin{column}{0.5\textwidth}
      \begin{block}{Flutter (Dart)}
        \begin{itemize}
          \item Собственный рендеринг UI
          \item Отличная скорость разработки (hot reload)
        \end{itemize}
      \end{block}
      \begin{block}{React Native (JS)}
        \begin{itemize}
          \item JS/TS стек, мост к нативным компонентам
        \end{itemize}
      \end{block}
    \end{column}
    \begin{column}{0.5\textwidth}
      \begin{block}{Kotlin Multiplatform (KMP)}
        \begin{itemize}
          \item Общая бизнес-логика, нативные UI
          \item Поддержка Авроры :-)
        \end{itemize}
      \end{block}
      \begin{block}{.NET MAUI}
        \begin{itemize}
          \item C\#/.NET, единый стек Microsoft
        \end{itemize}
      \end{block}
    \end{column}
  \end{columns}
  \vspace{2mm}
  \textbf{Риски:} тонкости производительности, мосты к нативу, обновления экосистем
\end{frame}

\begin{frame}{Гибридные приложения и PWA}
  \begin{itemize}
    \item \textbf{Гибрид:} веб-страницы в нативном контейнере (WebView)
    \item \textbf{PWA:} офлайн-кеш, иконка на рабочем столе, уведомления (где поддерживается)
  \end{itemize}
  \begin{block}{Когда это уместно}
    \begin{itemize}
      \item Контентные приложения, MVP, быстрые пилоты
      \item Корпоративные внутренние инструменты
    \end{itemize}
  \end{block}
\end{frame}

\begin{frame}{Сравнение подходов (краткая таблица)}
  \small
  \begin{tabularx}{\linewidth}{l *{4}{>{\centering\arraybackslash}X}}
    \toprule
    & \textbf{Натив} & \textbf{Кросс-платф.} & \textbf{Гибрид} & \textbf{PWA} \\
    \midrule
    Производительность & \textbf{Высокая} & Средняя/Высокая & Средняя & Низк./Средн. \\
    Доступ к API       & \textbf{Полный} & Почти полный & Ограничен & Ограничен \\
    Стоимость          & Высокая & \textbf{Ниже} & Низкая & \textbf{Низкая} \\
    UX/Нативный вид    & \textbf{Лучший} & Хороший & Посредственный & Зависит от браузера \\
    \bottomrule
  \end{tabularx}
\end{frame}

% =====================================================
\section{Kotlin: обзор языка и примеры}
% =====================================================

\begin{frame}{Kotlin в двух словах}
  \begin{itemize}
    \item Современный, статически типизированный язык от JetBrains
    \item Официально поддерживается Google для Android
    \item Сильная совместимость с Java (JVM), есть KMP
    \item Философия: лаконичность, безопасность, практичность
    \item Поддерживает \textbf{функциональную / декларативную}, 
    \textbf{объектно-ориентированную}, \textbf{алгоритмическую} парадигмы, обобщенное программирование
    \item Библиотека реализует идеи реактивного программирования
  \end{itemize}
\end{frame}

\begin{frame}[fragile]{Hello, Kotlin}
\lstset{language=Kotlin}
\begin{lstlisting}
fun main() {
    println("Hello, Kotlin!")
}
\end{lstlisting}
\end{frame}

\begin{frame}[fragile]{Сравнение: Java vs Kotlin}
\begin{columns}[T,onlytextwidth]
\begin{column}{0.48\textwidth}
\textbf{Java}
\lstset{language=JavaLite}
\begin{lstlisting}
// Java
public class User {
  private final String name;
  public User(String name) { this.name = name; }
  public String getName() { return name; }
}
\end{lstlisting}
\end{column}
\begin{column}{0.48\textwidth}
\textbf{Kotlin}
\lstset{language=Kotlin}
\begin{lstlisting}
// Kotlin
data class User(val name: String)
\end{lstlisting}
\end{column}
\end{columns}
\end{frame}

\begin{frame}[fragile]{Null-safety}
\lstset{language=Kotlin}
\begin{lstlisting}
var name: String? = null     // Можно\ хранить\ null
println(name?.length)        // Безопасный вызов -> null
println(name ?: "нет имени") // Elvis-оператор
\end{lstlisting}
\begin{block}{Важно}
Оператор \texttt{!!} бросает NPE — используйте его осознанно.
\end{block}
\end{frame}

\begin{frame}[fragile]{Функции, лямбды, расширения}
\lstset{language=Kotlin}
\begin{lstlisting}
fun Int.isEven(): Boolean = this % 2 == 0

val nums = listOf(1, 2, 3, 4)
val evens = nums.filter { it.isEven() }  // [2, 4]
\end{lstlisting}
\end{frame}

\begin{frame}[fragile]{Сравнение с Java: обработка null и коллекции}
\begin{columns}[T,onlytextwidth]
\begin{column}{0.48\textwidth}
\textbf{Java}
\lstset{language=JavaLite}
\begin{lstlisting}
List<Integer> xs = Arrays.asList(1,2,3);
List<Integer> evens = xs.stream()
  .filter(x -> x % 2 == 0)
  .collect(Collectors.toList());
\end{lstlisting}
\end{column}
\begin{column}{0.48\textwidth}
\textbf{Kotlin}
\lstset{language=Kotlin}
\begin{lstlisting}
val xs = listOf(1, 2, 3)
val evens = xs.filter { it % 2 == 0 }
\end{lstlisting}
\end{column}
\end{columns}
\end{frame}

\begin{frame}{Где применяется Kotlin}
\begin{itemize}
  \item Android-разработка (официальный язык Google)
  \item Backend: Spring, Ktor, Micronaut
  \item Kotlin Multiplatform: общая логика для Android/iOS
  \item Kotlin/Native: десктопные, серверные, встраиваемые приложения
  \item Kotlin/JS: веб-разработка
  \item Скрипты и утилиты (JVM)
\end{itemize}
\end{frame}

\section{Основы Kotlin}

\begin{frame}[fragile]{Переменные: var и val}
\lstset{language=Kotlin}
\begin{lstlisting}
var x = 10     // изменяемая переменная
x = 20
val y = 30     // неизменяемая (аналог final)
\end{lstlisting}
\begin{block}{Важно}
Используйте \texttt{val} по умолчанию.  
\texttt{var} — только когда нужно изменять значение.
\end{block}
\end{frame}

\begin{frame}[fragile]{Базовые типы}
\begin{itemize}
  \item Числовые: \texttt{Int}, \texttt{Long}, \texttt{Double}, \texttt{Float}, \texttt{Short}, \texttt{Byte}
  \item Логический: \texttt{Boolean}
  \item Символьный: \texttt{Char}
  \item Строковый: \texttt{String}
\end{itemize}
\lstset{language=Kotlin}
\begin{lstlisting}
val a: Int = 42
val b: Double = 3.14
val c: String = "Kotlin"
\end{lstlisting}
\end{frame}

\begin{frame}[fragile]{if: базовый оператор}
\lstset{language=Kotlin}
\begin{lstlisting}
val a = 10
if (a > 0) {
    println("a положительное")
}
\end{lstlisting}
\end{frame}

\begin{frame}[fragile]{if-else}
\lstset{language=Kotlin}
\begin{lstlisting}
val a = -5
if (a >= 0) {
    println("Неотрицательное")
} else {
    println("Отрицательное")
}
\end{lstlisting}
\end{frame}

\begin{frame}[fragile]{if как выражение}
\lstset{language=Kotlin}
\begin{lstlisting}
val a = 5
val b = 7
val max = if (a > b) a else b
println("Максимум = $max")
\end{lstlisting}
\begin{block}{Важно}
\texttt{if} возвращает значение, как в функциональных языках.
\end{block}
\end{frame}

\begin{frame}[fragile]{if с несколькими ветками}
\lstset{language=Kotlin}
\begin{lstlisting}
val grade = 75
val result = if (grade >= 90) "Отлично"
             else if (grade >= 70) "Хорошо"
             else if (grade >= 50) "Удовлетворительно"
             else "Плохо"

println(result)
\end{lstlisting}
\end{frame}

\begin{frame}[fragile]{when: базовый вариант}
\lstset{language=Kotlin}
\begin{lstlisting}
val x = 2
when (x) {
    1 -> println("Один")
    2 -> println("Два")
    else -> println("Другое")
}
\end{lstlisting}
\end{frame}

\begin{frame}[fragile]{when без аргумента}
\lstset{language=Kotlin}
\begin{lstlisting}
val y = -3
when {
    y < 0 -> println("Отрицательное")
    y == 0 -> println("Ноль")
    else -> println("Положительное")
}
\end{lstlisting}
\end{frame}

\begin{frame}[fragile]{when с диапазонами и множеством значений}
\lstset{language=Kotlin}
\begin{lstlisting}
val x = 7
when (x) {
    0, 1 -> println("Ноль или один")
    in 2..5 -> println("От двух до пяти")
    !in 6..10 -> println("Не от шести до десяти")
    else -> println("От шести до десяти")
}
\end{lstlisting}
\end{frame}

\begin{frame}[fragile]{when и типы (smart cast)}
\lstset{language=Kotlin}
\begin{lstlisting}
fun process(obj: Any) = when (obj) {
    is String -> println("Строка длиной ${obj.length}")
    is Int -> println("Целое число: $obj")
    else -> println("Неизвестный тип")
}
\end{lstlisting}
\end{frame}

\begin{frame}[fragile]{when как выражение}
\lstset{language=Kotlin}
\begin{lstlisting}
val day = 3
val name = when (day) {
    1 -> "Понедельник"
    2 -> "Вторник"
    3 -> "Среда"
    else -> "Выходной"
}
println(name)
\end{lstlisting}
\end{frame}

\begin{frame}[fragile]{when с произвольными условиями}
\lstset{language=Kotlin}
\begin{lstlisting}
val a = 10
val b = 20
val relation = when {
    a < b -> "a меньше b"
    a > b -> "a больше b"
    else -> "a равно b"
}
println(relation)
\end{lstlisting}
\end{frame}

\begin{frame}[fragile]{Приведение типов}
\lstset{language=Kotlin}
\begin{lstlisting}
val n: Int = 42
val d: Double = n.toDouble()  // явное преобразование

val obj: Any = "строка"
if (obj is String) {
    println(obj.length)  // smart cast
}
\end{lstlisting}
\end{frame}

\begin{frame}[fragile]{Выражения в Kotlin}
\lstset{language=Kotlin}
\begin{lstlisting}
val result = try {
    10 / 2
} catch (e: Exception) {
    -1
}
println(result)
\end{lstlisting}
\begin{block}{Важно}
В Kotlin почти всё — выражения: \texttt{if}, \texttt{when}, \texttt{try}, лямбды.
\end{block}
\end{frame}

\begin{frame}[fragile]{Ввод и вывод}
\lstset{language=Kotlin}
\begin{lstlisting}
fun main() {
    println("Введите имя:")
    val name = readln()
    println("Привет, $name!")
}
\end{lstlisting}
\end{frame}

\begin{frame}[fragile]{Как надо писать практики}
\lstset{language=Kotlin}
\begin{lstlisting}
fun main() {
    println("Введите первое\ число:")
    val a = readlnOrNull()?.toIntOrNull()
    println("Введите второе\ число:")
    val b = readlnOrNull()?.toIntOrNull()

    if (a != null && b != null) {
        val sum: Int? = if ((b > 0 && a > Int.MAX_VALUE - b) ||
                            (b < 0 && a < Int.MIN_VALUE - b)) {
            null
        } else {
            a + b
        }

\end{lstlisting}
\end{frame}


\begin{frame}[fragile]{Как надо писать практики (продолжение)}
\lstset{language=Kotlin}
\begin{lstlisting}
        if (sum != null) {
            println("Сумма: $sum")
        } else {
            println("Ошибка: переполнение Int")
        }
    } else {
        println("Ошибка: введите корректные целые числа")
    }
}
\end{lstlisting}
\begin{block}{Особенности}
\begin{itemize}
  \item Проверка до сложения исключает переполнение
  \item Абсолютно безопасный ввод: \texttt{readlnOrNull() + toIntOrNull()}
\end{itemize}
\end{block}
\end{frame}



% =====================================================
\section{Резюме}
% =====================================================

\begin{frame}{Итоги лекции}
  \begin{itemize}
    \item Обсудили формат курса, оценивание и ожидания
    \item Поняли различия подходов к мобильной разработке и критерии выбора
    \item Познакомились с Kotlin: синтаксис, null-safety, примеры фунционального подхода
  \end{itemize} % Исправлена опечатка: </end{itemize> на \end{itemize}
\end{frame}


\end{document}